\documentclass[11pt]{article}
\usepackage[margin=2.5cm]{geometry}
\usepackage{amsmath,amssymb,amsthm,booktabs,longtable}
\newtheorem{theorem}{Theorem}
\newtheorem{lemma}[theorem]{Lemma}
\theoremstyle{remark}\newtheorem*{remark}{Remark}
\title{A four-inequality limit for the Part~II geometry\\ of the $7/8$ zero-free region}
\author{Tommi Niilola}
\date{October 16, 2026}
\begin{document}
\maketitle

\begin{abstract}
Part~II of the September~30, 2026 quasi-Riemann paper (OAI) proves a zero-free region
$\operatorname{Re}s>7/8$ for one fixed geometry $(b,l_x,\ell)=(1/8,17/48,1/6)$. We ask how far the
same exponent inequalities could go for any geometry. Four of them, combined with exact rational
multipliers, give $s_0\ge 4638705/5301659=0.8749535\ldots$, so the architecture has less than
$4.7\cdot10^{-5}$ of room below $7/8$. Assume the mean-square hypotheses $\mathrm{MS}(\kappa)$ for all
$\kappa\in[\kappa_0,1]$, with the uniformity of OAI's (17.87). Under them the same method gives exact lower bounds that match the numerical optima to about $10^{-5}$, and a
$\kappa_0$-free bound $s_0\ge 167/192=7/8-1/192$. Both bounds survive dropping the condition
$M+\ell\le1$ once its row-norm cost is accounted for. A smaller floor bin gains at most $1/320$. The statements concern this exponent
architecture only. They are not barriers for other methods.
\end{abstract}

\section{The inequalities}
Line numbers refer to \texttt{paper.tex} of OAI. Write $l_y=l_x+b$, $M=l_x+l_y$, $h=1-l_x+\ell$,
$z_0=17/50$, and let $s_0$ be the target boundary ($\sigma_0$, $\beta_*$ in OAI). An independent reader
derived each inequality below from the source before comparing it with our transcription. Each one
matches the paper or is weaker than it, so it is a necessary condition for any geometry for which the
argument closes.
\begin{itemize}
\item[(N)] $b,l_x,\ell\ge0$ and $l_x\ge\ell$. The paper requires $P_a\ge1$ (line 8344) and
$l_x-d\ge3/16$ for $d\in[0,\ell]$ (line 8594).
\item[(S)] $M+\ell\le1$. The paper fixes $M+\ell=1$ (line 6873).
\item[(L)] Low estimate (Prop.~probe-low, lines 8564--8642). The bound
$l_x/2+b/12+\max_{d\in[0,\ell]}(-d+\mathrm{loss}(d)/2)$ must not exceed
$C(s_0)=l_x/2+s_0-1+h/6$. Taking $d=0$ and dropping $\mathrm{loss}(0)\ge0$ gives
\[ s_0\ \ge\ \tfrac56+\tfrac{b}{12}+\tfrac{l_x-\ell}{6}. \]
\item[(H)] High exponent (eq.~stage-high-exponent, line 6092). At row exponent $d=h$, which is the worst
case on $[1/2,h]$ (line 16133), with frequency bin $\delta$ and $g=x\delta\ell$:
\[ a-s_0+h(z_0-\tfrac16)-a\,l_y-\tfrac{\ell}{2}+x\delta\ell+h\bigl(R+\tfrac{\delta}{2}-z_0\bigr)\le0,
\qquad a=\tfrac{1+\delta}{2}, \]
for every $\delta\in[1/50,\kappa]$, $\kappa=2s_0-1$, and $x\in[0,1/2]$. Here
$R=R^*(\delta,x,\ell/h)$ is the row-count exponent of Prop.~19.2 with the long-range cost of
Lemma~17.6, an infimum over the cutoff $t$ of a maximum over the length $r$. We use (H) only at
$x=1/2$. Lemma~\ref{lem:count} evaluates $R^*$ for that case, and it does not cover other $x$.
\item[(F)] Floor bin $\delta_0=1/50$ with $R=1$ (eq.~floor-bound, line 15845): (H) with
$\delta=1/50$, $x=1/2$, $R=1$.
\end{itemize}
In (H) and (F) the terms in $z_0$ cancel, so the bounds below do not depend on the choice
$z_0=17/50$.

\paragraph{Relaxation below $7/8$.} Several lemmas of OAI are stated for $\sigma_0\ge7/8$. We do not
use any of their conclusions. We only use inequalities (L), (H), (F), each of which says that an
exponent the argument must make negative is negative, together with the count of Prop.~19.2 at
capacity $c=1/(6\max(3/4,\kappa))=2/9$, the largest capacity Lemma~plain allows. A hypothetical
extension of the proof to $s_0<7/8$ needs at least these conditions, and granting it the
unproved lemmas only enlarges its feasible set. So the infeasibility below is a valid limit for the
architecture.

\section{The count lemma}
\begin{lemma}\label{lem:count}
Let $2/3\le\kappa_0\le1$, $0<\delta<5\kappa_0/6$, $x=1/2$, $q=\delta/2$, $c=2/9$. For
$t\in[1,3/2]$ and $r\in[t-1/2,t]$, $m=t-r$, put
$f=e(r)-\delta r$ for $r\ge1$, with $e(r)=1$ for $r\le1/\kappa_0$ and $(1+5\kappa_0 r)/6$ otherwise, and
$f=\min\bigl(1-\delta r-2q z_i,\ 1-2\delta m-2q z_p\bigr)$ for $r<1$, with $z_i=\min((1-r)/2,\sigma)$,
$z_p=\min(c(1-2m),\sigma)$, $\sigma$ the supply cap. Then
$R^*(\delta,\sigma)=\inf_t\sup_r f\ \ge\ R_\infty(\delta)$, where
\[ R_\infty(\delta)=1-\frac{\delta(16+14t^*)}{37},\qquad
t^*=\frac{5/6-16\delta/37}{5\kappa_0/6-23\delta/37}, \]
if $1/\kappa_0\le t^*\le3/2$, and $R_\infty(\delta)=1-\delta$ otherwise.
\end{lemma}
\begin{proof}
Lowering the caps raises $f$, so $R^*(\delta,\sigma)\ge R^*(\delta,\infty)$. Without caps,
$z_i=(1-r)/2$ and $z_p=\frac29(1-2m)$. For $r<1$ the two entries of the minimum are
$1-\delta/2-\delta r/2$, decreasing in $r$, and $1-2\delta/9-14\delta m/9$, increasing in $r$. They
meet at $r=(28t-5)/37\in[t-1/2,1)$, so $\sup_{r<1}f=S(t)=1-\delta(16+14t)/37$, a decreasing function
with $S(3/2)=1-\delta$. For $r\ge1$, $\sup f=\max\bigl(1-\delta,\,L(t)\bigr)$, where
$L(t)=1/6+t(5\kappa_0/6-\delta)$ is increasing and present only for $t>1/\kappa_0$. Since
$L(1/\kappa_0)=1-\delta/\kappa_0<S(1/\kappa_0)$, the infimum over $t$ of
$\max(1-\delta,S,L)$ is $S(t^*)=L(t^*)$ when $t^*\le3/2$, and $1-\delta$ otherwise.
\end{proof}
At $\kappa_0=1$ this gives $R_\infty(19/50)=55889/82850$ and $R_\infty(2/5)=2129/3245$. The values
agree with the reduced-model code to within its grid error ($6\cdot10^{-5}$).

\section{The bound at $\kappa_0=1$}
\begin{theorem}\label{thm:main}
Let $\kappa_0=1$. Every geometry $(b,l_x,\ell)$ that satisfies (N), (S), (L) and (H) has
\[ s_0\ \ge\ \frac{4638705}{5301659}=0.8749535\ldots. \]
\end{theorem}
\begin{proof}
If $s_0<5/6$, (L) and (N) give $s_0\ge5/6$, a contradiction. Let $s_0\ge5/6$. Then
$\kappa\ge2/3$, so $\delta=19/50$ and $\delta=2/5$ are admissible in (H). With
Lemma~\ref{lem:count} and $R\ge R_\infty$, (H) becomes
\begin{align*}
(\mathrm H_1)\quad s_0&\ge \tfrac{172483}{124275}(1-l_x)+\tfrac{48208}{124275}\ell-\tfrac{69}{100}b,\\
(\mathrm H_2)\quad s_0&\ge \tfrac{13526}{9735}(1-l_x)+\tfrac{3791}{9735}\ell-\tfrac{7}{10}b.
\end{align*}
Write each inequality as $c_b b+c_l l_x+c_e\ell+c_s s_0\le \rho$ and add them with the multipliers in
Table~\ref{tab:mult}. The coefficients of $b$, $l_x$ and $\ell$ cancel exactly and the coefficient of $s_0$
is $-1$, which gives the bound.
\end{proof}
\begin{table}[h]\centering\small\setlength{\tabcolsep}{3.5pt}
\begin{tabular}{lrrrrrr}\toprule
Inequality & multiplier $\lambda$ & $c_b$ & $c_l$ & $c_e$ & $c_s$ & $\rho$\\\midrule
(L) & $4307228/5301659$ & $1/12$ & $1/6$ & $-1/6$ & $-1$ & $-5/6$\\
(S) & $331477/5301659$ & $1$ & $2$ & $1$ & $0$ & $1$\\
(H$_1$) & $1706710/15904977$ & $-69/100$ & $-172483/124275$ & $48208/124275$ & $-1$ & $-172483/124275$\\
(H$_2$) & $116053/1445907$ & $-7/10$ & $-13526/9735$ & $3791/9735$ & $-1$ & $-13526/9735$\\\midrule
$\sum\lambda\,\cdot$ & & $0$ & $0$ & $0$ & $-1$ & $-4638705/5301659$\\\bottomrule
\end{tabular}
\caption{Exact multipliers for Theorem~\ref{thm:main} (\texttt{symproofk.py} checks every entry).}
\label{tab:mult}
\end{table}

The obstruction is a three-way balance. (L) rewards large $l_x$ and small $\ell$, (H) charges for
small $1-l_x$ through the count $R$, and (S) ties the three lengths together. The floor, the
intermediate range and the length condition $l_y-d-11b/6\ge1/12$ do not enter. The audit marked the
last two as stronger than the paper needs, so it matters that the proof does not use them.

\section{Conditional variants}
Assume $\mathrm{MS}(\kappa)$ for every $\kappa\in[\kappa_0,1]$, with the uniformity over twists, symbol
signs, rowwise profiles and heights required in OAI's (17.87). A single estimate at $\kappa_0$ is not
enough. Then $e(r)=1$ for $1\le r\le1/\kappa_0$, and the copy argument of Lemma~17.6, run from
$\mathrm{MS}(\kappa_0)$, gives $e(r)\le(1+5\kappa_0r)/6$ beyond. This changes only the long-range cost in
Lemma~\ref{lem:count}. The same method gives the following exact bounds. Each certificate has four
rows, chosen from (L), (S), (H) at $x=1/2$, (F) and $b\ge0$. Appendix~\ref{app:cert} lists their
multipliers.
\begin{table}[h]\centering
\begin{tabular}{lllll}\toprule
$\kappa_0$ & rows used & exact bound & decimal & numerical optimum\\\midrule
$1$ & L, S, H($19/50$), H($2/5$) & $4638705/5301659$ & $0.8749535$ & $0.87497$\\
$97/100$ & L, S, H($19/50$), H($2/5$) & $49696303/56869303$ & $0.8738687$ & $0.87386$\\
$19/20$ & L, S, H($19/50$), H($9/25$) & $7753639/8880447$ & $0.8731136$ & $0.87312$\\
$9/10$ & L, S, H($17/50$), $b\ge0$ & $336893/386718$ & $0.8711593$ & $0.87117$\\
$4/5$ & L, S, F, $b\ge0$ & $167/192$ & $0.8697917$ & $0.8698$\\\bottomrule
\end{tabular}
\caption{Lower bounds under $\mathrm{MS}(\kappa)$ for all $\kappa\in[\kappa_0,1]$. The numerical optima are from the reduced model,
without the intermediate convenience bound $76/75-2\delta/3$.}
\end{table}

\begin{theorem}
For every $\kappa_0$, the inequalities (N), (S), (L) and (F) force $s_0\ge167/192=7/8-1/192$.
\end{theorem}
\begin{proof}
(F) reads $s_0\ge\frac{203}{150}(1-l_x)+\frac{53}{150}\ell-\frac{51}{100}b$. Adding (L), (S), (F)
and $-b\le0$ with multipliers $103/128$, $25/384$, $25/128$ and $25/768$ cancels $b$, $l_x$ and $\ell$
and gives $s_0\ge167/192$. None of these rows involves the count, so the bound holds for any
mean-square input.
\end{proof}

\begin{remark}
No mean-square improvement can move this architecture below $7/8-1/192$. Going further would need a
change in the floor bin $\delta_0=1/50$, which is tied to the Euler region $x_r\ge51/100$, in the
low estimate, or in the relation $M+\ell=1$.
\end{remark}

\section{Robustness}
\paragraph{The condition $M+\ell\le1$ is not needed.} OAI uses $M+\ell=1$ once, in the row-norm
lemma (line 8246, ``Since $M'+\ell'-1=-3d$''). Put $\sigma=M+\ell-1$. Without the condition the
available dual length satisfies $T_d\le H-3d+\sigma$. Retained dual dyads have
$v+\ell_b\le y\le T_d$, so $\max(H,v+\ell_b)\le H+(\sigma-3d)_+$, and the row-norm exponent grows by at
most $(\sigma-3d)_+$. The loss formula $M'+(1+3\ell'-2M')/4$ comes from the definition of $T_d$ and does
not involve $\sigma$. After Cauchy--Schwarz, the low estimate at $d=0$ becomes
\[ (\mathrm L_\sigma)\qquad s_0\ \ge\ \tfrac56+\tfrac{b}{12}+\tfrac{l_x-\ell}{6}+\tfrac12(b+2l_x+\ell-1)_+ . \]
Replacing (S) by the linear branch of (L$_\sigma$) gives the same bounds. At $\kappa_0=1$ the
multipliers of (L), (L$_\sigma$), (H$_1$), (H$_2$) are $3644274/5301659$, $662954/5301659$,
$1706710/15904977$, $116053/1445907$, with bound $4638705/5301659$. For the floor bound they are
$259/384$, $25/192$, $25/128$, $25/768$, with bound $167/192$. A geometry with $M+\ell>1$ would help only if the penalty
slope $1/2$ fell below about $0.077$.

\paragraph{The floor bin.} With a general floor bin $\delta_0=2a-1$, the rows (L), (S), (F) and
$b\ge0$ give exactly
\[ s_0\ \ge\ \frac23+\frac{1+2\delta_0}{5+6\delta_0}, \]
which is $167/192$ at $\delta_0=1/50$ and tends to $13/15$. Lowering $a$ below $51/100$ costs nothing
in the local correction (lines 4144--4170). The good-prime exponent $4-6x_r-6z_r$ is $-1.07$ at
$x_r=0.505$, $z_r=17/50$, and the ramified exponents are at most $1-2x_r-\epsilon_0$. In the full
system, though, the floor binds only when $\kappa_0\lesssim0.8$. Numerically, the lower bound at
$\kappa_0=1$ and $0.9$ is unchanged for $\delta_0=0.01$ and $0.004$, and at $\kappa_0=0.8$,
$\delta_0=0.004$ the count binds again at $0.86813$.

\paragraph{A bilinear low estimate.} If the low bound $l_x/2+b/12$ improved to $l_x/2+b/12-\eta l_x$, the
bound would fall steadily. About $\eta=0.07$, a saving of roughly $Z^{-0.03}$ over Cauchy--Schwarz,
would reach $0.85$, and no other constraint binds before $\eta=0.2$. However, the low-separated sum
$Q^{-1/2}\sum_m w_mA_mB_m$ equals the probe $I_\eta$ itself (eq.~low-separated, line 3470). An
unconditional saving over Cauchy--Schwarz would itself be the stronger low estimate, so it cannot
be supplied as a separate input. Measured cancellation at finite scales cannot establish the
uniform saving either. It reflects the zeros present at those scales: without zeros off the line
the ratio is about $Z^{C(1/2)-l_x/2-b/12}\approx Z^{-0.36}$. This does not rule out a new
unconditional bilinear argument, but such an argument would amount to a new way of dividing the
probe between the low and high ranges.

\section{Other approaches tested}
We also tested several modifications of the architecture. Each item below records the specific
obstruction we met and says whether it is an exact argument or a numerical finding. These are
limits of the approaches tested, not impossibility results for other methods.
\begin{itemize}
\item \emph{Mean-square input, floor bin, $M+\ell\le1$} (exact; Sections 4 and 5). Mean-square
input alone stops at $167/192$. A smaller floor bin gains at most $1/320$, and only for
$\kappa_0\lesssim0.8$. Relaxing $M+\ell\le1$ is paid back by the row norm.
\item \emph{Cancellation across floor rows} (analytic). After Poisson summation, the common-contour
version leads to an off-diagonal problem at row range $H\asymp D$, where the available generic
estimates give no saving. This does not exclude estimates that use more arithmetic structure.
\item \emph{Decomposition by top row dyads} (analytic). The filter to sixth-power-free rows
costs $Z^{h/6}$, which is what the decomposition gains.
\item \emph{Free ceiling on the Poisson frequencies, and savings in the floor-row count}
(numerical, reduced model, not certified). Neither gains anything at $\kappa_0=1$
($0.87498$ against $0.87497$), and the gain is below $0.002$ at $\kappa_0=0.8$.
\item \emph{A cubic probe} (exact at the exponent level). The cubic pole at $z=1/3$ fits the
signal budget, and the generic cubic Gram cost ($P_a^{1/3}$ in place of $P_a^{1/6}$) is small.
However, the generic cubic large sieve controls the fourth moment of Lemma~plain only up to
length $M/2$, short of the length $M$ the argument needs, with a deficit of $m/3$ coming from the
cube copies. Moreover, assuming the uncentred threshold is $(1-1/k)M$, the comparison and
centred branches of its induction need a comparison length
$M/k\le L<(\frac12-\frac1k)M$. Such an $L$ exists only for $k>4$. A cubic probe therefore needs a
different fourth-moment argument, not a transplant of the sextic one.
\item \emph{A bilinear low estimate} (Section~5). The sum is the probe itself, so a saving would
have to be an unconditional estimate, which amounts to a new way of dividing the probe.
\end{itemize}
Within the audited exponent model, none of the modifications tested moves the boundary
materially. Going beyond this architecture appears to need a new signal mechanism or a
different low-separation estimate.

\paragraph{Reproducibility.} The script \texttt{symproofk.py} checks every statement in exact
rational arithmetic. The script \texttt{certs\_cond.py} re-verifies all seven certificates. It writes
them, with every row coefficient, to the machine-readable file \texttt{cond\_certificates.json}.
The checker \texttt{certcheck.py}
independently verifies the $\kappa_0=1$ result for $s_0\le437/500$ with a 1013-box certificate.

\appendix
\section{Exact multipliers of the remaining certificates}\label{app:cert}
Each block lists the four rows, their multipliers, and for (H) the value $R_\infty(\delta)$ from
Lemma~\ref{lem:count} at that $\kappa_0$. With $R=R_\infty$, the row coefficients follow from the
formulas in Section~1. In every block the coefficients of $b$, $l_x$ and $\ell$ cancel exactly and that of
$s_0$ is $-1$. Table~\ref{tab:mult} gives the case $\kappa_0=1$.
{\small
\begin{longtable}{llll}\toprule
& row & multiplier $\lambda$ & $R_\infty(\delta)$\\\midrule
% generated by certs_cond.py (tab_cond.tex)
\multicolumn{4}{l}{$\kappa_0=97/100$: $s_0\ge49696303/56869303$}\\
& (L) & $46109803/56869303$ & \\
& (S) & $3586500/56869303$ & \\
& (H) $\delta=19/50$ & $30799925/170607909$ & $211961/317525$\\
& (H) $\delta=2/5$ & $1478575/170607909$ & $1151/1775$\\
\addlinespace
\multicolumn{4}{l}{$\kappa_0=19/20$: $s_0\ge7753639/8880447$}\\
& (L) & $2396745/2960149$ & \\
& (S) & $563404/8880447$ & \\
& (H) $\delta=19/50$ & $1324609/8880447$ & $10749/16225$\\
& (H) $\delta=9/25$ & $365603/8880447$ & $30697/45025$\\
\addlinespace
\multicolumn{4}{l}{$\kappa_0=9/10$: $s_0\ge336893/386718$}\\
& (L) & $207987/257812$ & \\
& (S) & $49825/773436$ & \\
& (H) $\delta=17/50$ & $49825/257812$ & $102997/149475$\\
& $b\ge0$ & $3347/1546872$ & \\
\addlinespace
\multicolumn{4}{l}{$\kappa_0=4/5$: $s_0\ge167/192$}\\
& (L) & $103/128$ & \\
& (S) & $25/384$ & \\
& (F) & $25/128$ & \\
& $b\ge0$ & $25/768$ & \\
\addlinespace
\multicolumn{4}{l}{$\kappa_0=1$, without (S): $s_0\ge4638705/5301659$}\\
& (L) & $3644274/5301659$ & \\
& (L$_\sigma$) & $662954/5301659$ & \\
& (H) $\delta=19/50$ & $1706710/15904977$ & $55889/82850$\\
& (H) $\delta=2/5$ & $116053/1445907$ & $2129/3245$\\
\addlinespace
\multicolumn{4}{l}{$\kappa_0=4/5$, without (S): $s_0\ge167/192$}\\
& (L) & $259/384$ & \\
& (L$_\sigma$) & $25/192$ & \\
& (F) & $25/128$ & \\
& $b\ge0$ & $25/768$ & \\
\bottomrule
\end{longtable}}
\end{document}
